\documentclass{article}
\usepackage[UTF8]{ctex}
\usepackage{amsmath,amssymb,amsfonts}
\usepackage{enumitem}
\usepackage{geometry}
\geometry{a4paper, margin=1in}
% === 错题本排版优化（不参与题目标识） ===
\setlength{\parindent}{0pt}
\setlength{\parskip}{0.35em plus 0.10em minus 0.08em}
\setlength{\emergencystretch}{3em}
\setlist{itemsep=0.25em, topsep=0.40em, parsep=0pt}
\allowdisplaybreaks[2]
% === 错题本排版优化结束 ===
\begin{document}
\section*{湖南大学2024年高等代数真题}
\section*{132．湖南大学2024年高等代数试卷}

\subsection{题目1}
1．设 \(f(x)=(x+1)^{2024}-x^{2024}-1\) ．

\begin{itemize}
  \item[（1）] 求 \(f^{\prime}(x)\) ．
  \item[（2）] 求 \(f^{\prime}(x)\) 的所有复数根及在复数域和实数上的不可约因式分解．
  \item[（3）] 判断 \(f(x)\) 是否有重根，并说明理由．
\end{itemize}

\subsection{题目2}
2．判断题．正确的请简要证明，错误的请举出反例．

\begin{itemize}
  \item[（1）] 已知 \(V=W_{1} \oplus W_{2}\) ，则对任意的 \(\alpha \in V\) ，有 \(\alpha \in W_{1}\) 或 \(\alpha \in W_{2}\) 。
  \item[（2）] 多项式 \(p(x)\) 在数域 \(K\) 上不可约，则 \(p\left(x^{2}\right)\) 在数域 \(K\) 上也不可约。
  \item[（3）] \(n\) 为偶数，则存在 \(A, B \in M_{n \times n}(\mathbb{R})\) ，满足对任意的 \(0 \neq \alpha \in \mathbb{R}^{n}\) ，都有 \(A \alpha, B \alpha\) 线性无关。
\end{itemize}

\subsection{题目3}
3．设 \(n\) 阶矩阵 \(A=\left(\begin{array}{rrrrr}1 & -1 & & & \\ & 1 & -1 & & \\ & & & \\ & & & 1 & -1 \\ & & & & 1\end{array}\right)\) ．求 \(A^{-1}\) ．

\subsection{题目4}
4．记 \(N(A)=\{\lambda \in \mathbf{C} \mid \lambda A\) 和 \(A\) 相似 \(\}\) 。

\begin{itemize}
  \item[（1）] \(A=\left(\begin{array}{lll}0 & 0 & 1 \\ 0 & 0 & 0 \\ 0 & 0 & 0\end{array}\right)\) ，求 \(N(A)\) ．
  \item[（2）] \(A\) 不是幂零矩阵，证明：\(N(A)\) 为有限集．
\end{itemize}

\subsection{题目5}
5．已知 \(V\) 为有限维线性空间， \(\mathcal{A}\) 为 \(V\) 上的线性变换．

\begin{itemize}
  \item[（1）] 证明： \(\operatorname{dim} V=\operatorname{dim} \operatorname{Ker} \mathcal{A}+\operatorname{dim} \operatorname{Im} \mathcal{A}\) ．
  \item[（2）] 证明： \(\mathcal{A}\) 可逆的充要条件是 \(\mathcal{A}\) 为单射．
  \item[（3）] 举例说明 \(V\) 为无限维线性空间时，（2）不成立．
\end{itemize}

\subsection{题目6}
6．设 \(A\) 是数域 \(K\) 上的 \(n\) 阶矩阵，证明 \(r(A)=r\) 的充要条件是：存在两个线性无关的向量组

\[
\alpha_{1}, \alpha_{2}, \cdots, \alpha_{r} \in K^{n}, \beta_{1}, \beta_{2}, \cdots, \beta_{r} \in K^{n} .
\]

使得

\[
A=\alpha_{1} \beta_{1}^{\mathrm{T}}+\alpha_{2} \beta_{2}^{\mathrm{T}}+\cdots+\alpha_{r} \beta_{r}^{\mathrm{T}} .
\]

\subsection{题目7}
7．设 \(A\) 为复数域上的 \(n\) 阶可逆矩阵，\(A^{2}\) 在复数域上可相似对角化，证明：\(A\) 在复数域上可相似对角化。

\subsection{题目8}
8．设 \(A=\left(a_{i j}\right)\) 为3阶实正定对称矩阵，且 \(a_{i j} \in\{-1,0,1\}\) ，求矩阵 \(A\) ，并证明你的结论。
\end{document}
